\chapter{Discussion and Conclusions}\label{ch:conclusions}

\section{Findings}\label{sec:conclusions-findings}\label{sec:conclusions-amp}

\paragraph{Derived exactly.} For the stationary AR(1) chain under
coordinatewise OU corruption, the joint score is $\score(x,t) = -\Qt x$
with $\Qt = (e^{-2t}\Sigz + \Deltat I)^{-1}$, equal by the
score--posterior identity to $(e^{-t}\E[a|x] - x)/\Deltat$ (RQ1). Its
coupling structure has a closed-form lifecycle: exactly tridiagonal at
$t = 0$, filling one diagonal per order in $t$ by the band-fill law
\eqref{eq:g-bandfill}, and returning to the identity at rate $e^{-2t}$.
In the bulk, the single quantity $q(\alpha, t) \in (0, |\alpha|]$ governs
both the posterior correlation length and the geometric decay of evidence
influence, with limits $q \to 0$ as $t \to 0$ and $q \to |\alpha|$ as
$t \to \infty$. The posterior of the noisy chain lives on a tree-shaped
factor graph, on which the two BP sweeps compute all smoothing marginals
in $O(K)$ message updates; on the Gaussian chain the messages close
algebraically in the Gaussian family, the recursion is the Kalman filter
followed by the RTS smoother, and the cost becomes $O(K)$ (RQ2). The
error of a radius-$r$ local score estimator decays exactly as $q^r$
(RQ4). On the same model, the per-node AMP/TAP closure returns the exact
score (the mean fixed point is closure-independent) but a different
variance, with existence boundary, breakdown time $t_c(\alpha)$, and
critical coupling $\alpha_c = \sqrt{2} - 1$ in closed form
(Appendix~\ref{app:amp}); all of these formulas were re-verified
independently.

\paragraph{Measured numerically.} The two Gaussian score routes agree to
$10^{-15}$--$10^{-14}$ across the tested $(K, \alpha, t)$ grid, with an
independently written implementation reproducing the agreement. On
Laplace chains, Gaussian-projected BP measured against validated grid
quadrature has median relative score error $0.39$ at $t = 0.02$,
decaying below $10^{-3}$ for $t \gtrsim 0.9$; the error is amplified at
small $t$ by the exact factor $e^{-t}/\Deltat$, is worst for bimodal
innovation laws, and decreases with stronger chain correlation, while
the score direction remains accurate (median cosine $\geq 0.92$
throughout) (RQ3). At equal locality budget, truncating the inference is
more accurate than truncating the score matrix at small $t$ ($12.3\%$
against $21.0\%$ at $t = 0.05$).

\paragraph{The Gaussian-to-Laplace comparison.} The structural lesson is
where exactness is lost: the tree, the two-sweep schedule, and the
score--posterior identity survive any innovation law; what breaks is the
representability of messages in a finite-dimensional family. The
practical lesson is when it matters: the Gaussian closure is essentially
free at moderate and large diffusion times, and expensive precisely in
the small-$t$ regime where the score magnitude is amplified, though even
there the score direction is preserved much better than its magnitude.

\section{What the model suggests but does not establish}

Within the Gaussian benchmark, local score computations are near-optimal
away from intermediate diffusion times, the required receptive field
scales as $\xi \log(1/\varepsilon)$ with $\xi = 1/\log(1/q)$, and a
local algorithm is preferable to a banded linear readout. These are
properties of exact inference in one controlled model. Whether trained
score networks on real sequence data inherit them is untested here and
would require learning experiments.

\section{Limitations}

(i) The exact solutions concern one-dimensional state per frame and
linear AR(1) dynamics. (ii) The non-Gaussian results are numerical: grid
BP is a reference only within its validated resolution regime
($\sqrt{2t} \gtrsim 3\,\dd x$), and the Laplace closure error is a
measured curve, not a formula. (iii) The locality laws are proved for
the Gaussian bulk; their non-Gaussian counterparts are only structurally
delimited. (iv) No learning is performed. (v) An earlier conjecture on
low-rank structure of the Laplace $K = 2$ posterior Hessian failed
outside a narrow regime and is not asserted. (vi) How score errors
accumulate along the reverse generation trajectory was not studied.

\section{Future work}

Three directions follow directly from what was measured. First, mixture
messages: the single-Gaussian closure fails worst on bimodal beliefs, so
a two-component mixture closure on the same sweeps would separate the
representation error from the model mismatch, and is the clearest next
algorithmic step. Second, discrete chains: with a finite alphabet the BP
messages are finite vectors and inference is exact with no closure at
all, giving a second solvable benchmark. Third, reverse dynamics:
integrating the reverse SDE under exact, Gaussian-projected, and
truncated scores would show where along the trajectory the closure and
locality errors matter for generation, turning the pointwise error
measurements of this thesis into dynamical ones.
